
\section{怎样写比较复杂的议论文}

《中学语文教学大纲》对高中二年级学生开始提出“着重培养比较复杂的议论能力”的作文训练的要求。这个“比较复杂”也是跟“比较简单”相对而言的。哪些地方“比较复杂”呢？一般说来，比较复杂的议论文所要论述的问题比较复杂，难于一下弄清；论证的层次比较多，一个中心论点往往要绕几个分论点，各个分论点又有完整的论证过程；论证的方法多种多样，立论和驳论要结合运用，特别要求有严密的逻辑性；文章的结构方式比较多，需要综合运用，等等。那么，我们在写比较复杂的议论文时，应该注意哪些问题呢？

1. 从分析题目类型、弄清题意出发，明确而突出地提出一个恰当的中心论点。中心论点的确立和提出，是写好比较复杂的议论文的重要一步。要确立和提出一个恰当的中心论点，首先应从分析题目入手。

比较复杂的议论文的题目，大致有如下几种类型：

（1）题目就是中心论点。这一类最好办，照题论证就易，如《无私才能无畏》、《进步从知不足开始》、《知识就是力量》、《天才来自勤奋》、《文贵创新》，等等。

（2）题目只提出要论说的向题或对象。如《先天下之忧而忧，后天下之乐而乐》、《青春啊，应该怎样度过？》、《说“傻子”》、《从“谢谢”所想到的》、《“羡慕”小议》，等等。这就需要分析这个问题或对象的内在含义，明确论说的重点，从而提出自己正确、鲜明的主要见解，即中心论点。

（3）题目限定要从事物的对照中发议论。如《想和干》、《骄与馁》、《荣辱辨》、《“竞争”面面观》、《嫉妒、羡慕和奋起直追》，等等。这类题目，要求我们从议论对象的互相联系和对照比较中提炼出自己的看法，确立文章的中心论点。

（4）题目提供进行批驳或规劝的问题或对象。如《知识越多越反动吗？》、《考试并没有结束》、《切莫蹉跎岁月》等。这类题目的中心论点实际已经点明，就在句意的反面。它们与用肯定语气的判断句表示的题目不同的地方，就在于它们应写成以驳论为主的议论文，而不能以正面立论为主。

（5）题目本身只提供议论材料，观点比较隐晦，需要贯通题目的内容及其内在联系与含义，引申提炼出一个有启发性的论题，从而进一步确立自己的中心论点。如《昨天、今天、明天》、《马路、山路和人生之路》《笨鸟先飞》等；或者只提供一则寓言（象高考语文试卷的作文题，曾提供过刘基《郁离子·捕鼠》和韩非《韩非子·说林上》中关于“毁树容易种树难”的译文，叫写读后感）、一幅漫画（象高考语文试卷的作文题，曾提供过一幅《这下面没有水，再换个地方挖》的漫画，叫写一段说明性文字和一篇议论文）。我们就要从它们的象征意义、比喻意义、拟人意义或含蓄其中的其他意义进行发掘，确立我们的中心论点。

前四种类型的题目如何审题立意，我们已在《怎样进行构思》一章中作了说明，第五种类型的题目如何审题立意，我们将放在下一章《怎样写读后感、观后感》中去谈。

应该注意：尽管有些议论文题目可以从不同的角度确立几个中心论点，但在一篇不长的议论文中只能提出一个中心论点，而且这个中心论点一旦确定之后，就要贯彻始终，行文中不能转移中心论点，否则就叫犯了“论点中途转移”的毛病。

2.围绕中心论点，确立几个分论点。为了清楚、透彻地证明中心论点的正确性，往往要从几个方面去论述。如果每个方面有一个观点，每一个观点又需要加以证明的话，那么在每个观点论证的范围内，每个观点就成了论点，就全文而言，这样的论点叫做分论点。比如课文《民族的科学的大众的文化》，其中心论点是“民族的科学的大众的文化，就是人民大众反帝反封建的文化，就是新民主主义文化，就是中华民族的文化。”为了证明这个中心论点是正确的，毛泽东同志用了三个分论点，即“这种新民主主义文化是民族的”、“这种新民主主义的文化是科学的”、“这种新民主主义的文化是大众的，因而即是民主的”，进行了充分、有力的论证。可以看出，在一篇比较复杂的议论文中，中心论点是“纲”，分论点是“目”。中心论点要统帅各个分论点，而各个分论点则必须对中心论点起到证明、补充或发挥的作用。只有这样，文章才能条分缕析，纲举目张。

应该注意：各个分论点要紧紧围绕着一个中心论点，不允许分论点说的是与中心论点无关或扣得不紧的话；各个分论点之间必须要互相关连，一环扣着一环，一层深入一层，做到清楚而又透彻地证明中心论点，而不可杂乱无章，没有联系；各个分论点也必须有充分、可靠、典型的事实或理论作为论据，而不可一笔带过，因为如果不能证明分论点确是确凿可信的话，中心论点的正确性也就发生动摇。

3. 要综合运用多种立论方法，使文章具有严密的逻辑性。从论证的方式看，议论文一般可分立论和驳论两种。关于立论的基本方法，我们在《怎样写比较简单的议论文》一章中已谈到了例证、引证、比较、类比和分析五种。在写比较复杂的议论文中，人们还经常运用归纳和演绎这两种立论方法。

归纳是从个别到一般的推理方法，即从许多个别典型事例的分析中，归纳出一个共同的一般的结论。把这种推理方法用到立论中来，就要以许多确凿的事实材料作论据，把自己的观点、主张树立起来，做到先分说、后总说。例如，我们具体分析了个别事实，象西汉杰出的历史家司马迁“读万卷书，行万里路”，不仅在书本中学，而且到实践中学，以毕生的精力写出了巨著《史记》；晋代杰出的书法家王羲之临池学书，池水尽黑，终于使他的草隶冠绝古今；唐代大诗人白居易年轻时刻苦读书，以致“口舌成疮，手肘成胝”，终于成为唐代三大顶峰诗人之一；宋代的范仲淹少小无依，立志苦学，终于成为一代杰出的政治家和文学家；明代学者宋濂，由于家穷，向人借书来学，甚至冒着严寒到百里外求师；清代大思想家顾炎武一生刻苦读书，勤作笔记，终于编写成著名的《日知录》，新文化革命运动的主将鲁迅，连别人喝咖啡的时间都不放过，用来读书写作，成为伟大的革命家、思想家和文学家；人民音乐家冼星海即使穷困得住在上海的亭子间时，仍然把身子伸出阁楼的天窗之外坚持拉小提琴；大画家徐悲鸿在交不起房租、吃不上饭的情况下，无视一切，坚持作画；当代的苏阿芒刻苦自学，成为著名的世界语诗人；苏元星发奋学习天文，连新婚之夜都不忘观察星星象，终于发现了新星……我们从这些个别事实中，就可以归纳出“成就来自勤奋好学”这样一个科学论断。这种推理方法是以具体的事实作为根据来证明论点的，而这些事实又是客观存在的。因此这种推理方法有严密的逻辑性，说服力很强。

演绎是从一般到个别的推理方法。这种方法，是根据已知的一般原理或结论，来推断出关于个别事物的结论的。把这种推理方法运用到立论中来，就要先摆出已为人们所理解和接受了的科学论断，再依据这个科学论断，对所提出的事实或问题进行推论分析，得出新的符合客观规律的论断，做到先总说，后分说。例如，我们在提出“成就来自勤奋好学”这个科学论断后，接着就可以推论出古代的“诗圣”杜甫，大散文家韩愈、欧阳修，大小说家曹雪芹等，近代的文学家郭沫若、茅盾、巴金、老舍，政治家毛泽东、周恩来、陈毅，艺术家梅兰芳、常香玉、夏菊花等等在各个领域取得重大成就者，无不是勤奋好学的人。演绎推理是以科学的理论作根据进行推理，根据是可靠的，逻辑力量也是很强的，具有说服力。

归纳和演绎，两者互相联系，互相补充。在写作一篇比较复杂的议论文时，可以分开运用，更多的是把二者结合起来，互相配合，交替使用，形成“总说——分说——总说”的结构形式。

4.综合运用立论和驳论，掌握各种驳论方法。除了立论之外，驳论也是议论文的重要论证方式。立论是为了建立正确的论点，重在“立”，驳论是为了批驳错误的观点，重在“破”。但是，“破”和“立”是辩证统一的，因此，立论和驳论也不能截然分开，常常是互相联系、交叉使用的。例如，为了写好立论的文章，有时也需要提出错误的观点予以批驳。这样，有了对立面的存在，文章就能写得有生气，正面论点也可以申述得更充分。同样，在写作以批驳为主的驳论文时，也需要在批驳错误论点的基础上，简要地阐明正确的论点。因为从根本上说，批驳错误论点的目的正是从反面来证明正确的观点。

驳论的方法主要有：

（1）反驳论点，这是驳论中常用的方法。我们可以用事实证明对方论点的错误；可以直接剖析错误论点及其危害性，可以建立起与对方论点相对立的新论点，针锋相对地驳倒对方，也可以把对方的错误论点加以合理的引申，最后暴露它的谬误与荒唐，从而驳倒对方。例如课文《魔鬼的笛音》，针对莫斯科霸权主义者 讲 谤 中国企图“征服整个地球”、“加剧国际紧张局势”、“威胁整个世界”的反动论点，运用大量无可辩驳的事实进行驳斥，指出了谎言同事实、谬论同真理之间的矛盾，从而彻底驳倒了这种无耻的诬蔑。

（2）反驳论据。错误的论点，往往建立在虚假的论据之上，揭露了论据的虚伪性，对方的论点也就不攻自破，无法成立了。例如鲁迅在《“友邦惊诧”论》中，就采用了这种方法。文章最后引用的《申报》南京专电中的两则消息，是铁的事实，它有力地驳斥了国民党反动派对学生的诬蔑，学生既然不是破坏“社会秩序”的罪魁，那么国民党反动派建立在这种诬蔑之上的论据则是一派谎言，从而使敌人的反动论点“友邦人士，莫名惊诧，长此以往，国将不国”根本不能成立了。

（3）反驳论证。这种驳论方法，着眼于揭露对方在论证中的逻辑错误，证明对方的论点和论据之间没有必然的逻辑关系，这样，对方的论点也就不能成立了。比如某老师明明是南方人，但有的同学却认为他是北方人，理由是“他经常吃面条，而且能讲一口准确流利的普通话”。我们只要指出这二者之间没有必然的逻辑联系，那么“他是北方人”这个论点就是站不住的。

由于反驳论据、反驳论证，归根到底也是为了反驳论点，这三者之间有着不可分割的关系，所以在驳论文的写作中常常结合起来运用。

应该注意：在运用这些方法写作以驳论为主的议论文时，一定要根据反驳对象的不同，错误性质的不同而采取不同的态度。对于人民内部的错误观点、模糊认识，要热情帮助，耐心说服，切忌无限上纲，说过头话；批驳敌人的谬论，则应义正词严，毫不留情，但也要做到有理有据，不能辱骂和恐吓。

5.要根据实际情况，灵活运用各种结构方式。议论文一般都需要“提出一个什么问题，接着加以分析，然后综合起来，指明问题的性质，给以解决的方法。”（毛泽东《反对党八股》）我们应该掌握议论文这种结构方式的基本顺序，打好写作议论文的基础。但是，文章的结构千变万化，贵在创新，我们还得根据内容和特定的思路去安排文章的结构，不要总是一成不变地写成三大段，形成公式化，那就不好了。

议论文的开头，往往在提出问题的同时，就点明中心论点，或者提出批驳对象，树起靶子。一般要开门见山，立即接触论题，切忌兜圈子，说套话。

议论文的主要部分在于分析问题，组织各个有关的论据来证明论点或反驳谬论，要求条理清楚，层次分明。各个层次之间的关系，可以采取从不同角度排出并列的分论点来证明中心论点的“平列式”，如毛泽东同志的《民族的科学的大众的文化》、《反对党八股》等，也可以采取由远到近、从古到今、由小到大、由表及里地排列层层深入的各个分论点的“推进式”，如刘少奇同志的《个人和集体》、陶铸同志的《崇高的理想》等，也可以采取通过对事物本身各个发展阶段的纵的对比，或者与另一事物的横的对比，来排列分论点的“比较式”，如邓小平同志的《讲实事实求是》、胡雪同志的《惊人相似的一幕》等；还可以采取先从反面论述，再从正面论述，或者先从正面论述，再反面论述的“正反”式，如贾谊《论积贮疏》，就是先从反面论述不积贮的危害性，再从正面论述积贮的大有利，在古今正反两方面对照中，论证了积贮为治国之本这个中心论点；还可以采取针对靶子逐点批驳、引出论点的“批驳”式，如鲁迅的《“丧家的”“资本家的乏走狗”》，先驳梁实秋“说我是资本家的走狗，是那一个资本家，还是所有的资本家？”，指明“这正是‘资本家的走狗’的活写真”；再驳梁“我还不知道我的主子是谁”，指明是“丧家的”“资本家的走狗”；再用事实驳梁“这一套本领，我可怎么能知道呢？”，指明梁“给主子嗅出匪类（‘学匪’）”、“以济其‘文艺批评’之穷”的阶级实质，所以还得在“走狗”之上，加上一个形容字：“乏”。

议论文的结尾是全文发展的自然结果。写法很多，或总结全文，以强调中心论点，或展示未来，以鼓舞斗志；或抒发情怀，以增强文章的感染力；或照应前文，以使首尾呼应……总之，结尾应深化中心论点，语言要简练有力。既不要草草了事，使人感到文章似乎还没有写完，也不要画蛇添足，絮絮叨叨地再写些套话、大话之类的东西。
\newpage